% Source: https://artofproblemsolving.com/community/c3h1771925_mock_sprint_27__tricky_sequence
% Stable identity: aops-topic-1771925-q01
Call a number of the form $p^{2^j}$ ($p$ prime, $j\ge0$) a \emph{special} number: $2,3,4,5,7,9,11,13,16,\dots$. Writing each exponent in a prime factorisation in binary shows that every integer $n>1$ is a product of distinct special numbers in exactly one way (for example $2^5\cdot3=2\cdot16\cdot3$).

\textbf{Claim: after $1$, Xavier writes exactly the special numbers in increasing order.} Suppose the terms so far are $1$ and all special numbers less than $m$. The products of distinct terms are exactly the numbers whose special factors are all less than $m$. A number not of this kind has a special factor $q\ge m$, so it is at least $q$, and the smallest such number is therefore the smallest special number $\ge m$. So the next term is the next special number, which proves the claim by induction.

The composite special numbers up to $1000$ are $p^2$ for the primes $p\le31$ ($11$ numbers: $4,9,25,49,121,169,289,361,529,841,961$), $p^4$ for $p=2,3,5$ ($16,81,625$) and $2^8=256$; $2^{16}$ and $3^8$ are too large. The first term above $1000$ is the prime $1009$ (no square, fourth power or eighth power lies in $1001$ to $1008$). So Xavier writes $11+3+1=\boxed{15}$ composite numbers.
